\ifmostrarGabarito
\begin{minted}{python}
from typing import List

def calcular_media_movel(
    precos_diarios: List[float], 
    tamanho_janela: int
) -> List[float]:
    if len(precos_diarios) < tamanho_janela or tamanho_janela <= 0:
        return []
        
    medias = []
    n_janelas = len(precos_diarios) - tamanho_janela + 1
    
    for i in range(n_janelas):
        janela = precos_diarios[i : i + tamanho_janela]
        media_janela = round(sum(janela) / tamanho_janela, 2)
        medias.append(media_janela)
        
    return medias
\end{minted}

\textbf{Explicação:}
A função valida primeiramente se o número de preços é suficiente para formar ao menos uma janela de tamanho $K$. Caso positivo, calcula a quantidade de intervalos deslizantes possíveis ($N - K + 1$) e percorre os dados com fatiamento de listas (\texttt{precos\_diarios[i : i + tamanho\_janela]}). A média de cada janela é calculada com \texttt{sum()} e \texttt{round()}, acumulando os resultados na lista final.

\vspace{0.3cm}
\textbf{Critérios de Correção e Penalidades Sugeridas (Erros Comuns):}
\begin{itemize}[leftmargin=*, noitemsep, topsep=2pt]
    \item \textbf{Não validar tamanho mínimo ou janela nula/negativa} (\textbf{--4,0 pontos}): Não testar se \code{len(precos\_diarios) < tamanho\_janela} ou se \code{tamanho\_janela <= 0}, gerando exceção ou divisão por zero.
    \item \textbf{Erro no fatiamento da janela} (\textbf{--6,0 pontos}): Usar \code{precos\_diarios[i : tamanho\_janela]} em vez de \code{[i : i + tamanho\_janela]}, mantendo o ponto final fixo em vez de deslizar a janela.
    \item \textbf{Erro na quantidade de iterações (\emph{off-by-one})} (\textbf{--5,0 pontos}): Iterar apenas até \code{len(precos) - tamanho\_janela}, deixando de calcular a última média móvel.
    \item \textbf{Cálculo incorreto da média} (\textbf{--5,0 pontos}): Dividir pelo tamanho total da lista em vez de \code{tamanho\_janela}, ou somar elementos incorretos.
    \item \textbf{Arredondamento e formato de retorno} (\textbf{--5,0 pontos}): Omitir \code{round(..., 2)} nas médias calculadas ou retornar apenas a última média em vez da lista completa.
\end{itemize}
\else
\linhasResposta{17}
\clearpage
\linhasResposta[FOLHA DE CONTINUACAO / RASCUNHO -- QUESTAO Q10]{30}
\clearpage
\fi
